% Independent supporting proof fragment. No standalone preamble.
\section{Matched positive-transport certificates and block refinement}
\label{supp:block-transport-review}

I retain the actual rounded midpoint-quantile nodes throughout this
comparison. Let $\nu(dy)=\phi(y)\mathbf1_{[-L,L]}dy$, let its mass be
$a$, and divide its CDF range into $N$ cells $I_c=[a_c,b_c]$ of mass
$a/N$. The node $c_c$ is the actual dyadic rounded node, with
$a_c\leq c_c\leq b_c$ established by an outward interval certificate.
The finite measure $\tau=(a/N)\sum_c\delta_{c_c}$ has the same raw
mass $a$ as $\nu$. The grid, retained region and mass are identical
for both certificates below.

\subsection{A competent first-order comparator}

For a continuously differentiable function $f$, the fundamental theorem
of calculus on each side of a node gives
\[
 \int_{I_c}\{f(y)-f(c_c)\}\phi(y)dy
 =\int_{a_c}^{c_c} f'(z)\{\Phi(a_c)-\Phi(z)\}dz
  +\int_{c_c}^{b_c} f'(z)\{\Phi(b_c)-\Phi(z)\}dz.
\]
The exchange of integrals is justified on the two compact integration
triangles. Consequently
\[
 |(\nu-\tau)(f)|\leq\int|f'(z)|\rho_0(dz),
\quad
 \frac{d\rho_{0,c}}{dz}=
 \begin{cases}
 \Phi(z)-\Phi(a_c),&a_c\leq z<c_c,\\
 \Phi(b_c)-\Phi(z),&c_c<z\leq b_c.
 \end{cases}
\]
The density is nonnegative. Its value at the single node is irrelevant.
The total mass is $\sum_c\int_{I_c}|y-c_c|\phi(y)dy$. Thus this
comparator preserves local transport distances and the same dispersion
moments as the second-order method. It is more informative than a
whole-box supremum times a maximal coordinate displacement.

For $v\neq0$, integration by parts gives the cell transform
\[
 \mathcal M_{\rho_{0,c}}(v)=\frac{
 e^{vc_c}\{2\Phi(c_c)-\Phi(a_c)-\Phi(b_c)\}
 +\mathcal M_{\nu|[c_c,b_c]}(v)
 -\mathcal M_{\nu|[a_c,c_c]}(v)}{v},
\]
where
$\mathcal M_{\nu|[u,t]}(v)=e^{v^2/2}\{\Phi(t-v)-\Phi(u-v)\}$.
At zero its exact value is
\[
 c_c\{2\Phi(c_c)-\Phi(a_c)-\Phi(b_c)\}
 +2\phi(c_c)-\phi(a_c)-\phi(b_c).
\]
The sign of the denominator is retained for negative rates. For tiny
nonzero rates, positivity and support provide
$e^{-|v|L}\mathcal M_{\rho_0}(0)\leq
 \mathcal M_{\rho_0}(v)\leq
 e^{|v|L}\mathcal M_{\rho_0}(0)$ instead of an unstable quotient.
The zero-rate formula is not substituted for a nonzero transform.

\subsection{The second-order certificate on the same law}

For $f\in C^2([-L,L])$, finite Taylor expansion about each actual node
gives
\[
 (\nu-\tau)(f)=\sum_c\mu_{1,c}f'(c_c)+\int f''d\rho_2,
\quad
 \mu_{1,c}=\int_{I_c}(y-c_c)\phi(y)dy.
\]
Here $\rho_2$ is positive, with density
$\int_{a_c}^z(z-y)\phi(y)dy$ below the node and
$\int_z^{b_c}(y-z)\phi(y)dy$ above it. In particular,
$\rho_{2,c}(I_c)=\frac12\int_{I_c}(y-c_c)^2\phi(y)dy$.
Choose rational $\bar\mu_{1,c}\geq|\mu_{1,c}|$ and put
$\rho_1=\sum_c\bar\mu_{1,c}\delta_{c_c}$. Then
\[
 |(\nu-\tau)(f)|\leq\int |f'|d\rho_1+\int |f''|d\rho_2.
\]
The absolute first-moment allowance does not assert favourable
cancellation between cells. For any block $D$ of complete consecutive
cells, define $\nu_D,\tau_D,\rho_{k,D}$ by restriction. Applying
the exact Taylor identity to $e^{vy}$ proves
\[
 \mathcal M_{\rho_{2,D}}(v)=
 \frac{\mathcal M_{\nu_D}(v)-\mathcal M_{\tau_D}(v)
   -v\sum_{c\subset D}\mu_{1,c}e^{vc_c}}{v^2}
 \quad(v\neq0).
\]
At zero the half-second-moment formula is used. Positivity, support and
outward arithmetic remain necessary when cancellation widens this
quotient. A rigorous support enclosure is an available fallback.

Both certificates follow from established Taylor transport. The second
cell identity is a specialisation of Cerone (2002, Theorem~3,
equations~(2.13)--(2.15), with $n=1$, $\alpha=a_c$, $\beta=b_c$,
$x=c_c$ and $w=\phi$; the endpoint choice is explicitly permitted in
the accompanying remark. I use the published identity as an antecedent,
rather than presenting the positive Peano measure as a new theorem.
The second-order certificate
uses additional smoothness and may be sharper for a particular
residual, but there is no uniform ordering of the two error allowances
merely because they share a rounded law. A matched comparison must
report both, using the same available derivative information and tail.

\subsection{Positive conditional moments and monotone refinement}

For any $p$ the finite-measure telescope is
\[
 \nu^{\otimes p}-\tau^{\otimes p}
 =\sum_{j=1}^p\nu^{\otimes(j-1)}\otimes(\nu-\tau)
                \otimes\tau^{\otimes(p-j)}.
\]
For a block partition $\mathcal P$ and $k\in\{0,1,2\}$, write
\[
 \Lambda_{j,k,D}=\nu^{\otimes(j-1)}\otimes\rho_{k,D}
                            \otimes\tau^{\otimes(p-j)}.
\]
Every measure is positive and finite. Its other coordinates contribute
their raw retained mass, including when their exponential rate is zero.
For the classically specified basket set
$X_i=e^{d_i+s_i\cdot y}$,
$T=\sum_iw_iX_i$, $U=T-1\geq0$,
$u=T_j$ and $v=T_{jj}\geq0$.
The previously derived residual derivative envelopes contain direct
nonnegative moments and Cauchy closures of $U|u|$ and $U^2|u|$.
Under any such block measure, put $\langle h\rangle_D=\int h\,d\Lambda_{j,k,D}$.
Then
\[
 \langle U|u|\rangle_D\leq
 \sqrt{\langle U^2\rangle_D\langle u^2\rangle_D},
\]
\[
 \langle U^2|u|\rangle_D\leq\min\left\{
 \sqrt{\langle U^2\rangle_D\langle U^2u^2\rangle_D},
 \sqrt{\langle U^4\rangle_D\langle u^2\rangle_D}\right\}.
\]
No probability normalisation or mass divisor occurs in these bounds.
If $D$ is the disjoint union of children $D_1,\ldots,D_m$, all
moments add exactly. Discrete Cauchy gives, for nonnegative $A_i,B_i$,
\[
 \sum_i\sqrt{A_iB_i}\leq
 \sqrt{\left(\sum_iA_i\right)\left(\sum_iB_i\right)}.
\]
For the two alternative closures, writing their child values $a_i,b_i$
gives
\[
 \sum_i\min(a_i,b_i)\leq
 \min\left(\sum_i a_i,\sum_i b_i\right)
 \leq\min(a_D,b_D).
\]
All coefficients in the residual derivative envelopes are nonnegative.
Direct moments are additive. Therefore each exact certificate envelope
is nonincreasing under a refinement of its transported-coordinate
blocks. Strict improvement can arise from failure of the relevant
Cauchy proportionality conditions or from a different preferred
alternative closure on different blocks.
Refining just this coordinate is sufficient; no $p$-dimensional block
Cartesian grid is required.

Raw outward numerical upper bounds need not display this exact
monotonicity, because independent intervals need not add tightly. For a
parent envelope with certified upper bound $b_D$ and children with
upper bounds $b_i$, the value
$\min\{b_D,\sum_i b_i\}$ is again a valid upper bound for the
parent integrand integral. Applying this operation recursively gives a
certified monotone numerical hierarchy. The same construction may be
used componentwise before summing positive envelope coefficients.
It does not justify intersecting unrelated interval quantities or
claiming that an interval midpoint is an upper bound.

\subsection{Exact factorisation, full tails and cost scope}

A required geometry polynomial has the form
$h=\sum_\alpha b_\alpha\prod_iX_i^{\alpha_i}$, with degree at
most four. For every block hybrid, independence gives
\[
 \langle h\rangle_D=\sum_\alpha b_\alpha
 e^{\sum_i\alpha_id_i}
 \mathcal M_{\rho_{k,D}}\!\left(\sum_i\alpha_is_{ij}\right)
 \prod_{\ell<j}\mathcal M_\nu\!\left(\sum_i\alpha_is_{i\ell}\right)
 \prod_{\ell>j}\mathcal M_\tau\!\left(\sum_i\alpha_is_{i\ell}\right).
\]
Signed polynomial coefficients are retained before taking outward
interval upper endpoints. The underlying polynomial's nonnegativity
permits intersecting an enclosure with $[0,\infty)$; it does not
permit changing a negative upper endpoint to zero. Positive rational
dominating first-moment weights are fixed across partitions so that
their aggregate measure is exactly additive.

All retained-grid transfer bounds require the same full continuous
Gaussian tail. For a second correction with global curvature bound
$M_2$, $|R_2|\leq M_2U^2/2$ yields
$\frac12M_2\mathbb E[U^2\mathbf1_{Y\notin[-L,L]^p}]$.
For the first correction $0\leq R_1\leq K e^{m_G+r\cdot Y+\lambda^2/2}U$,
so its corresponding tilted outside moment is used. For a cap the
minimum with $H\mathbb P(Y\notin[-L,L]^p)$ is valid. Gaussian
tilting evaluates each outside exponential monomial by
\[
 e^{b+\|s\|^2/2}
 \left[1-\prod_{j=1}^p\{\Phi(L-s_j)-\Phi(-L-s_j)\}\right].
\]
This tail is added once. Block refinement does not reduce, omit or
renormalise it. No price answer is an input.

Let $M\leq\binom{n+4}{4}$ be the number of distinct monomials and
let $B$ be the number of blocks. Rebuilding all $p$ factors for every
coordinate, block and monomial would require $O(Bp^2M)$ arithmetic
products, in addition to transform construction and polynomial
combination. The implementation caches the other-coordinate factors and
groups monomials by their transported-coordinate rate before evaluating
any partition. This requires $O(p^2M)$ product formation for a fixed
number of geometry polynomials, followed by $O(BpM)$ block dot products;
coincident rates reduce the latter count. Rational rate construction,
scalar-transform queries, factorial-series prefix preparation and
memory are additional costs. The retained operation counters identify
these costs rather than replacing them by a claim of linear dependence
on dimension. The finite one-
dimensional grid still has $N$ nodes. There is no $N^p$ price
enumeration in the certificate construction, nor a dimension-free
error or a quantum advantage theorem. The one-block envelope is the
aggregate certificate; the $N$-block endpoint is the competent
cellwise Cauchy certificate. Intermediate partitions provide a
measured construction-cost versus bound-tightness trade-off.

The refinement argument is an elementary conditioning consequence of
positive measures and Cauchy's inequality. The defensible research
question is whether its project-specific implementation supplies a
useful, evidenced pricing certificate trade-off against these matched
endpoints. The existence of a block hierarchy alone does not establish
a new general theory or literature priority.
